\section{CONCLUSION}

Focusing on the task of selenium nanoparticle (SeNP) cytotoxicity prediction, this study established an integrated research pipeline comprising benchmark dataset model comparison, dedicated selenium nanoparticle modeling, and automated screening of enhanced features. The findings demonstrate that on large-scale heterogeneous tabular nanoparticle datasets, both deep tabular architectures and tree-based models achieve effective predictive performance. However, in selenium nanoparticle scenarios characterized by limited sample sizes and a relatively homogeneous material scope, tree-based models exhibit superior stability and applicability.

Within the dedicated selenium nanoparticle dataset, 383 validated records were curated from 43 peer-reviewed publications \cite{chen_selenium_2008, li_reversal_2011, zheng_sialic_2011, liu_selenium_2012, wu_surface_2012, chen_surface_2012, yu_positive_2012, chen_peg-nanolized_2012, pi_pathway_2013, wu_induction_2013, feng_differential_2014, huang_selective_2013, li_functionalized_2013, zhang_enhancement_2013, jiang_gracilaria_2014, yu_ph-responsive_2014, zeng_construction_2015, liu_sequentially_2016, liu_facile_2016, yu_x-ray-responsive_2016, fu_rgd_2016, fang_targeting_2017, chen_efficient_2018, fang_dualtargeted_2018, huang_autophagy_2018, liu_highly_2018, wang_regulating_2016, hu_selenium_2019, zeng_potentiation_2019, huang_dual-targeting_2019, huang_highly_2019, zhang_systematic_2019, zheng_thermosensitive_2019, huang_chiralitydriven_2020, yang_lentinan-functionalized_2020, liu_rapid_2021, an_functionalized_2021, song_translational_2021, yang_atherosclerotic_2022, yang_functionalized_2021, liang_translational_2022, liu_specific_2022, du_highly_2022}, spanning the period from 2008 to 2022. While baseline physicochemical features provide preliminary predictive capabilities, the incorporation of enhanced variables—such as elemental Se content, long-term stability, Se 3d binding energy, crystalline phase encoding, and IC50—further enhances the fitting capacity of the models. Among these architectures, XGBoost yielded the highest regression performance under the enhanced feature configuration; nevertheless, because this configuration incorporates IC50, which possesses an intrinsic codependency with toxicological endpoints, this result is more appropriately designated as an auxiliary reference for mechanistic insight and validation. In contrast, Random Forest maintained a balanced compromise among precision, robustness, and interpretability, rendering it a more suitable comparative baseline for subsequent external validation and engineering implementation.

Feature selection results further indicate that the optimal combination of enhanced variables does not require the exhaustive set of candidate descriptors; utilizing solely IC50, elemental Se content, long-term stability, and Se 3d binding energy is sufficient to achieve predictive accuracy comparable to that of the complete feature set. Concurrently, this finding underscores the necessity of distinguishing IC50 from intrinsic material chemistry descriptors: the former operates primarily as auxiliary information coupled to toxicological responses, whereas the latter embodies transferable and cumulative feature engineering value within a single-material domain. Consequently, future data acquisition and experimental design should prioritize key descriptors that directly reflect the intrinsic chemical attributes of the nanomaterial.

Overall, this study delineates a structured pathway encompassing data curation, model benchmarking, and feature screening to support the development of selenium nanoparticle toxicity prediction platforms. Future work should integrate prospective experimental datasets and external validation cohorts to rigorously evaluate the generalization performance of models operating independently of response-correlated auxiliary variables across distinct batches and more complex application scenarios.

% 本文围绕纳米硒细胞毒性预测任务，构建了由基准数据集模型比较、硒纳米颗粒专用建模以及增强特征自动筛选组成的研究流程。研究结果表明，在大规模异质纳米颗粒表格数据上，深度表格模型与树模型均能够完成有效预测；而在样本规模有限且材料类型相对固定的硒纳米颗粒数据场景下，树模型表现出更高的稳定性和适用性。

% 在硒纳米颗粒专用数据集中，383 条有效记录来源于 43 篇文献\cite{chen_selenium_2008, li_reversal_2011, zheng_sialic_2011, liu_selenium_2012, wu_surface_2012, chen_surface_2012, yu_positive_2012, chen_peg-nanolized_2012, pi_pathway_2013, wu_induction_2013, feng_differential_2014, huang_selective_2013, li_functionalized_2013, zhang_enhancement_2013, jiang_gracilaria_2014, yu_ph-responsive_2014, zeng_construction_2015, liu_sequentially_2016, liu_facile_2016, yu_x-ray-responsive_2016, fu_rgd_2016, fang_targeting_2017, chen_efficient_2018, fang_dualtargeted_2018, huang_autophagy_2018, liu_highly_2018, wang_regulating_2016, hu_selenium_2019, zeng_potentiation_2019, huang_dual-targeting_2019, huang_highly_2019, zhang_systematic_2019, zheng_thermosensitive_2019, huang_chiralitydriven_2020, yang_lentinan-functionalized_2020, liu_rapid_2021, an_functionalized_2021, song_translational_2021, yang_atherosclerotic_2022, yang_functionalized_2021, liang_translational_2022, liu_specific_2022, du_highly_2022}，时间跨度为 2008--2022 年。基础特征能够提供初步预测能力，而引入元素 Se 含量、长期稳定性、结合能 Se3d、晶型编码及 IC50 等增强变量后，模型拟合能力进一步改善。其中，XGBoost 在含增强变量设置下取得最高回归结果；但由于该设置包含与毒理终点天然相关的 IC50，该结果更适合作为辅助分析和机理增强参考。随机森林则在精度、稳健性和可解释性之间保持了较好的平衡，更适合作为后续外部验证与工程实现中的比较基线。

% 特征筛选结果进一步说明，表现较优的增强变量组合并不需要依赖全部候选项，仅使用 IC50、元素 Se 含量、长期稳定性和结合能 Se3d 即可达到接近全特征集的预测精度。与此同时，这一结果也提示应将 IC50 与材料本征化学特征区分解释：前者更接近毒理响应相关辅助信息，后者则更能体现单一材料体系中可迁移、可积累的特征工程价值。因此，面向后续数据采集与实验设计，仍应优先关注能够反映材料本征化学属性的关键变量。

% 总体而言，本研究为纳米硒毒性预测平台建设提供了较为清晰的数据整理、模型比较和特征筛选路径。后续工作仍需结合新增实验数据和外部验证样本，进一步检验不依赖响应相关辅助变量时模型在独立批次和更复杂应用场景中的泛化能力。
